\documentclass[conference]{IEEEtran}

\newcommand{\todo}[1]{\textbf{TODO: #1}}
\begin{document}

\title{Vulture-AI: Context-Aware Vulnerability Prioritization with Retrieval-Augmented Generation}

\author{\IEEEauthorblockN{Anonymous Authors}
\IEEEauthorblockA{Anonymous Institute\\
Email: anonymous@example.com}}

\maketitle

\begin{abstract}
Vulnerability prioritization is challenging. In this paper, we introduce the Vulture Risk Score (VRS) and a retrieval-augmented generation (RAG) assistant...
\end{abstract}

\section{Introduction}
Vulnerabilities are typically ranked by CVSS \cite{first2023cvss}. However, real-world risk requires threat intelligence (e.g., EPSS \cite{jacobs2021epss}, CISA KEV \cite{cisa2023kev}) and asset context \cite{prioritization2022}.

\section{Related Work}
Recent work explores LLMs \cite{smith2023llm} and RAG \cite{rag2023sec} for security.

\section{Methodology}
The Vulture Risk Score (VRS) is defined as:
\begin{equation}
VRS = 100 \times (w_C C + w_E E + w_K K + w_N N + w_A A)
\end{equation}

\section{Experimental Setup}
We generated a benchmark of 40 CVEs, sourced from NVD \cite{nvd2023}. 

\section{Results}
Table \ref{tab:main} compares models.
\todo{insert result: explain NDCG@10 from main\_comparison.tex}
\begin{table}[h]
\centering
\begin{tabular}{cccc}
\hline
Model & NDCG@10 & NDCG@20 & Spearman \\
\hline
cvss\_only & 0.335 & 0.410 & 0.875 \\
epss\_only & 0.151 & 0.379 & 0.887 \\
combo\_baseline & 0.438 & 0.439 & 0.924 \\
vrs\_v1 & 0.434 & 0.435 & 0.924 \\
\hline
\end{tabular}
\end{table}

\section{Ablation \& Discussion}
Table \ref{tab:ablation} shows our ablation study.
\todo{insert result: explain the zero impact of KEV weight due to policy override}
\begin{table}[h]
\centering
\begin{tabular}{ccc}
\hline
Model & NDCG@10 & Spearman \\
\hline
vrs\_full & 0.434 & 0.924 \\
drop\_epss & 0.447 & 0.908 \\
drop\_kev & 0.434 & 0.924 \\
drop\_n & 0.440 & 0.918 \\
drop\_a & 0.512 & 0.932 \\
\hline
\end{tabular}
\end{table}

\section{Threats to Validity}
Synthetic benchmark constraints.

\section{Conclusion}
VRS outperforms CVSS.

\bibliographystyle{IEEEtran}
\bibliography{references}

\end{document}
